\section{Scheduling, tools, and execution authority}
\label{sec:boundary}

\subsection{Scheduling: when an update becomes visible}
\label{sec:boundary-scheduling}

Routing selects the next location; scheduling determines when its computation runs. The scheduler
also determines which record the computation reads and when later computations can read its update.

\begin{definition}[Superstep]
\label{def:superstep}
``Pregel computations consist of a sequence of iterations, called supersteps.'' By design, the model
``doesn't expose any mechanism for detecting order of execution within a superstep, and all
communication is from superstep $S$ to superstep $S + 1$''~\citep[\S1]{malewicz2010pregel}.
\end{definition}

A runtime using this model divides each superstep into three phases: i) \emph{planning} selects
computations whose inputs were written in the previous superstep; ii) \emph{execution} runs those
computations, whose writes remain invisible to other computations during this phase;
iii) \emph{update} applies the writes through the reducers of \cref{sec:state-reducers}. The next
superstep plans from the updated record. In the notation of \cref{sec:state-transition}, execution
computes $u_t$, and the update phase computes $\mathrm{merge}(s_t, u_t)$.

This sequence guarantees update order. The runtime applies a node's writes after the node returns,
then plans its successor. \textbf{A later node never reads a partly applied update.} Nodes in the
same superstep cannot see one another's writes. The runtime combines their writes in the update
phase. If two nodes write the same field, that field needs a reducer that combines their values;
otherwise, the runtime must reject the competing writes.

This guarantee applies only to the record. A tool can change a file during execution, before the
superstep ends. If the node then fails, the file remains changed, but the record lacks the update
that would describe the change. \textbf{Superstep visibility orders record updates. It does not
order external effects.} The tool boundary must control those effects.

\subsection{The tool contract}
\label{sec:boundary-contract}

The model receives a tool's name, description, and argument schema. The runtime holds its
implementation, enforces its operating constraints, and handles invalid arguments. The description
alone does not provide these controls.

\begin{definition}[Tool contract]
\label{def:tool-contract}
A \emph{tool contract} specifies five parts of a tool: i) its name and purpose, ii) typed arguments
and their constraints, iii) the typed success result and error representation, iv) its side-effect
class, and v) execution constraints, such as the directories it may access. The definition supplied
to the model covers the first two parts. The implementation and dispatcher determine the other
three, which the model cannot establish from the description.
\end{definition}

A method specification assigns obligations to the caller and the implementer
(\citealp[Specification structure]{mit6031-specs}).

\begin{definition}[Precondition]
\label{def:precondition}
``The precondition is an obligation on the client (the caller of the method). It is a condition over
the state in which the method is invoked.''
\end{definition}

\begin{definition}[Postcondition]
\label{def:postcondition}
``The postcondition is an obligation on the implementer of the method.'' ``In general, the
postcondition is a condition on the state of the program after the method is invoked, assuming the
precondition was true before.''
\end{definition}

For a tool, the precondition covers typed arguments and their constraints, including execution
constraints on the request. For example, a path may have to resolve inside the workspace. The
postcondition covers the result, error representation, and side effects. The precondition defines
when the tool's promises apply; the postcondition defines those promises.

A method specification may make no promises when the caller violates a precondition. An agent's
tool contract needs stricter rejection rules because the model can supply invalid arguments during
normal operation. The contract must state which violations the runtime detects, how it reports
refusal, and that a refused request has no side effects. \textbf{A tool contract must define how
the runtime handles violated preconditions.}

Our three tools illustrate this contract. \demoref{read\_file} takes a path and an optional line
range and returns text; \demoref{search} takes a regular expression and an optional glob and
returns matching lines with their locations; \demoref{edit} replaces one exact occurrence of a
string in a file. Only \demoref{edit} writes. All three take the workspace root as their first
parameter. Before opening a file, \demoref{read\_file} and \demoref{edit} resolve its path through
that root. \demoref{search} matches its glob against the root without resolving paths
(\cref{sec:boundary-demo}).

The function signature determines which arguments the dispatcher accepts, so it also determines
where configuration can be supplied. Tools read the run's configuration from the module. They do
not accept it as a parameter. The dispatcher binds model arguments against the signature, so
\textbf{a parameter the dispatcher has not already supplied can receive a value from the model}.
Adding a \demoref{config} parameter would let the model set it. The dispatcher itself supplies
the workspace-root parameter, which is absent from the model's schema. If a request names that
parameter anyway, binding fails because the dispatcher has already supplied the argument.

\subsection{The dispatch cycle}
\label{sec:boundary-dispatch}

The dispatch cycle has five operations, each of which can end the cycle.
\Cref{fig:dispatch-cycle} follows one request through them.

\begin{figure}[htbp]
  \centering
  \begin{adjustbox}{width=\linewidth}% Path policy may be checked inside the tool before the external effect.
% The dispatch cycle, top to bottom: one model request, through validation and policy, to an
% observation. The two dashed branches are refusals, and they rejoin at the same object the
% successful path produces, which is the point of the drawing: a refusal and a tool error are both
% observations, and neither is an exception.
\begin{tikzpicture}[node distance=4mm and 14mm]
  \node[tnode, align=center] (req) {tool request};
  \node[rnode, below=of req, text width=42mm] (val) {resolve the name; bind the arguments};
  \node[rnode, below=of val, text width=42mm] (pol) {resolve the path\\inside the workspace root};
  \node[rnode, below=of pol, text width=42mm] (exec) {perform the permitted operation};
  \node[rnode, below=8mm of exec, text width=42mm] (obs) {observation: \texttt{ok}, result, args};
  \node[tnode, align=center, below=of obs] (next) {state, trace event, next model input};
  \draw[flow] (req) -- (val);
  \draw[flow] (val) -- (pol);
  \draw[flow] (pol) -- (exec);
  \draw[flow] (exec) -- node[lab, right] {result, or tool error} (obs);
  \draw[flow] (obs) -- (next);
  \draw[back] (val.west) to[out=180, in=180, looseness=2.1]
    node[lab, left, pos=0.22, align=right] {unregistered, or arguments\\that do not bind} (obs.west);
  \draw[back] (pol.west) to[out=180, in=180, looseness=0.9]
    node[lab, left, pos=0.78, align=right] {outside the root} (obs.west);
\end{tikzpicture}
\end{adjustbox}
  \caption{The logical stages of one dispatch. Ellipses show the incoming request and the records
  that receive the observation; grey boxes show runtime work. Solid arrows follow successful
  checks. Dashed arrows show refusals, which also produce observations. A policy check may run
  inside the tool, as in our implementation, if it precedes the protected effect.}
  \label{fig:dispatch-cycle}
\end{figure}

\bi

\item \textbf{Resolve the name.} The dispatcher refuses requests for unregistered tools and lists
the available tools in its response. It builds the registry from the tool enumeration. A tool
absent from that registry cannot run, regardless of the definitions supplied to the model.

\item \textbf{Bind the arguments.} The dispatcher binds the model's arguments against the function's
signature before calling it. A wrong argument name therefore produces a readable refusal rather
than a \demoref{TypeError} from inside the tool. Binding checks names and arity. It does not check
types. The dispatcher coerces unambiguous type mismatches, avoiding an extra model call to correct,
for example, a string supplied where an integer was required.

\item \textbf{Check the policy.} Before opening a file, the tool resolves its path and refuses it
if it lies outside the workspace root. This check is the first operation in \demoref{read\_file}
and \demoref{edit}. The dispatcher converts any exception from this check into a refusal.

\item \textbf{Perform the permitted operation.} The function runs with the bound arguments and the
dispatcher's workspace root.

\item \textbf{Observe.} Success and failure alike become an \demoref{Observation} carrying the tool,
the arguments, an \demoref{ok} flag and a string result, and the dispatcher writes a
\demoref{tool\_result} event either way.

\ei

A refusal returns an observation to the model, allowing the reactive loop to continue. An
exception ends the run. If \demoref{edit} refuses an ambiguous target, the model can use the
observation to retry with more context. Raising an exception would end the run at that first
imprecise edit. Exceptions are reserved for unrecoverable failures, such as an unwritable trace.

The tools node dispatches multiple calls from one response sequentially within one superstep.
Each call receives an observation, including failed calls. Omitting a response to any call makes
the next model message malformed.

\subsection{Evidence of a request and its execution}
\label{sec:boundary-artifacts}

One tool call appears in six places, each of which supports different claims.
\Cref{tab:boundary-artifacts} follows the first \demoref{read\_file} call of the planning agent's accepted
run, \demoref{bc284c6b}, through these records and processing stages.

\begin{table}[htbp]
  \centering
  \begin{adjustbox}{width=\linewidth}\begin{tabularx}{\linewidth}{@{}L{0.25\linewidth}L{0.36\linewidth}X@{}}
    \toprule
    \textbf{Artifact} & \textbf{What it contains} & \textbf{What it establishes} \\
    \midrule
    Model input & The tool definitions and the conversation so far. The trace does not record
      it. & Which operations and descriptions were supplied to the model. \\
    Model response, and the \demoref{model\_call} event (step 3) & A tool call with an identifier, a
      name and arguments. The event keeps only the name, \demoref{read\_file}. & That the model
      requested an operation. It does not establish execution. \\
    \demoref{tool\_request} event (step 5) & The tool name and the arguments as the dispatcher
      received them. & That the request reached the dispatcher and began the cycle. \\
    Dispatcher & Registry lookup, argument binding, path resolution. No event of its own; a failed
      check would appear as the reason in the result event. & That the request satisfied the
      contract, or why it failed a check. \\
    \demoref{tool\_result} event (step 6) & \demoref{ok=true}, 1,665 characters, which is the length
      of the file. & That the function ran and returned, and how much text came back. \\
    State and next input & One record appended to \demoref{observations}, one tool message appended
      to \demoref{messages}. & What the next model call, at step 7, was allowed to see. \\
    \bottomrule
  \end{tabularx}\end{adjustbox}
  \caption{Records and processing stages for one \demoref{read\_file} call on
  \demoref{testcode/BenchmarkTest00283.py}, from run \demoref{bc284c6b}, steps 3 to 7. Step numbers
  count trace events rather than supersteps. The message history retains the provider's raw
  request; the trace records the request received by the dispatcher.}
  \label{tab:boundary-artifacts}
\end{table}

Step 4, between the response and the request, is the routing event: decision \demoref{tools},
\demoref{by="model"}, because the reply carried a tool call (\cref{sec:control-routing}). It records
the selected control location. It does not describe tool execution.

The model response and the \demoref{tool\_result} event support different claims. A tool call in
a model response establishes what the model requested. It does not establish a file-system
operation. If the model produces both a request and a plausible result in a transcript, neither
establishes execution. In this example, the runtime's \demoref{tool\_result} event has
\demoref{ok=true}, establishing that the function ran and returned. A result event can also report
a refusal before function entry. The successful event does not establish the returned text's
contents, because it stores only the result's length.

\subsection{Execution authority}
\label{sec:boundary-authority}

The model proposes actions; the runtime executes or refuses them. Each responsibility belongs to
an identifiable component on the request's execution path. Establishing that the checks cover
\emph{every} execution path requires a separate argument.

\begin{definition}[Complete mediation]
\label{def:complete-mediation}
``To attempt to secure this resource, server needs to check all accesses to the resource.
`Complete mediation.' '' ``Server will put a `guard' in place to mediate every request for this
particular resource. Only way to access the resource is to use the guard.''~\citep[\S5]{mit6033-lec20}
\end{definition}

For an agent, a resource is anything a tool can affect, including files, processes, or the network.
The request is a tool call, and the guard consists of the checks performed before its effect.
Complete mediation requires coverage of all execution routes. A correct check on one tested route
says nothing about a route that bypasses it. We must therefore enumerate the routes to each
protected operation. Coverage is most easily lost through an overlooked route, such as a new tool
or an argument that identifies a file without being named a path. Complete mediation
ensures that the runtime consults the policy. It does not establish that the policy is correct.

The term ``control'' involves two separate distinctions.

\bi

\item \textbf{Decision timing and ownership.} Some choices are fixed when we write the program and
others are computed during execution. Timing alone does not identify who controls the choice.
A routing function and a limit check both run during execution, but both apply rules the
programmer wrote.

\item \textbf{Proposal and execution.} A requested action, a claimed result, or a declaration that
the work is finished takes effect only through the runtime's execution path. A model's statement
that it ran the tests does not establish that tests executed.

\ei

\begin{table}[htbp]
  \centering
  \begin{adjustbox}{width=\linewidth}\begin{tabularx}{\linewidth}{@{}XL{0.5\linewidth}@{}}
    \toprule
    \textbf{Guarantee} & \textbf{Where it is enforced} \\
    \midrule
    Only registered tools run & The registry built from the tool enumeration, consulted before
      anything is called. \\
    Arguments bind to the tool's signature & Signature binding and coercion in the dispatcher,
      before the call. \\
    No file outside the workspace is written & \demoref{workspace.resolve}, called by
      \demoref{edit}, the one tool that writes, on its path argument. \\
    No file outside the workspace is read & Not enforced. \demoref{read\_file} resolves its path,
      but the glob of \demoref{search} never reaches \demoref{workspace.resolve}. \\
    The pinned fixture is not modified & Every tool works on a disposable copy of the fixture, and
      the diff compares that copy with the pinned fixture. \\
    An exhausted budget prevents another model call & The check each model node makes before it
      issues its call. \\
    A plan is revised at most once per run & The tools node, which records a stall for revision
      only while \demoref{replanned\_at} is unset, and the routing function after tools, which reads
      that record (\cref{tab:control-backedges}). \\
    \bottomrule
  \end{tabularx}\end{adjustbox}
  \caption{Desired guarantees and the components that enforce them. The implementation does not
  enforce one of these guarantees. Prompt instructions enforce none of them.}
  \label{tab:boundary-authority}
\end{table}

\demoref{workspace.resolve} fully resolves both the root and the candidate path before comparing
them. Resolution collapses \demoref{..} segments and follows symlinks, preventing either from
escaping the root. A string-prefix comparison would not provide the same check:
\demoref{/tmp/ws} is also a prefix of the unrelated directory \demoref{/tmp/ws-evil}. Joining an
absolute path to the root yields the absolute path itself. The check refuses it if its resolved
destination lies outside the root.

\subsection{The demo: scheduling and request refusal}
\label{sec:boundary-demo}

Our planning agent runs on LangGraph's runtime, which organizes execution ``into multiple steps, following
the Pregel Algorithm/Bulk Synchronous Parallel model'', and whose execution phase is specified so
that ``channel updates are invisible to actors until the next step''~\citep{langgraph-pregel}. Our
graph has no parallel branches, so each superstep runs one node. In run \demoref{bc284c6b}, steps
3 and 4 belong to the controller's superstep: first the model call, then the routing event. The
routing function reads the record after the controller's reply has been merged. Steps 5 and 6
belong to the next superstep, when the tools node reads that reply and dispatches its call. The
dispatcher writes both trace events during execution. The observation enters
\demoref{observations} only after the node returns, at the update boundary.

No recorded run contains a refused request. We therefore construct this demonstration by sending
requests directly to the dispatcher, outside a graph, using a fresh workspace for the worked
alert. We report the dispatcher's responses. The first request calls \demoref{read\_file} on
\demoref{testcode/BenchmarkTest00283.py}. The dispatcher permits it and returns the same 1,665
characters as step 6 of the recorded run. The second request changes only the argument.

\bi

\item \textbf{The request.} \demoref{\{"name": "read\_file", "args": \{"path":
"../../etc/passwd"\}\}}, which no model produced.

\item \textbf{The refusing check.} \demoref{workspace.resolve}, the first operation of
\demoref{read\_file}, raises \demoref{PathEscapeException}, and the dispatcher catches it.

\item \textbf{The record.} An observation with \demoref{ok=false} whose result is the refusal,
\demoref{refused: \textquotesingle../../etc/passwd\textquotesingle{} resolves outside the workspace root}, and a
\demoref{tool\_result} event carrying the same reason.

\item \textbf{Evidence that no file was opened.} The trace holds a \demoref{tool\_request} event and a
failed \demoref{tool\_result} event with a reason and no character count. That shows the request was
refused. The code establishes that no file was opened: the check precedes the first file operation.
The trace does not record file operations and cannot establish this on its own.

\ei

The two requests use the same tool and dispatcher. Their arguments determine which execution
branch the runtime selects. This identifies the enforcing component described in
\cref{sec:boundary-authority}.

The same exercise exposes a route that bypasses the check. The dispatcher permits a
\demoref{search} whose glob reaches outside the workspace. With the glob
\demoref{../*-withheld/testcode/*.py} it returns lines from the benchmark cases the workspace
withholds in a sibling directory. With enough \demoref{../} segments, it returns the matching line
of \demoref{/etc/passwd}. Both results have \demoref{ok=true}.
\textbf{The path check guards \demoref{read\_file} and \demoref{edit}; it does not guard
\demoref{search}, whose glob never passes through \demoref{workspace.resolve}.} Our boundary against
writing outside the workspace is complete because \demoref{edit} is the only tool that writes.
The implementation does not completely mediate reads. The repair is to resolve every file
\demoref{search} matches through \demoref{workspace.resolve} before reading it.

\subsection{Reading}

\S1 and \S2 of \citet{malewicz2010pregel}, for supersteps and the rule that
communication crosses superstep boundaries in one direction only; the section ``Specification
structure'' of \citet{mit6031-specs}, for preconditions and postconditions; and \S5 and \S6 of
\citet{mit6033-lec20}, the guard model and the ways complete mediation is bypassed.

For reference while writing HW1: \citet{langgraph-pregel} for the plan, execution and update
phases as our runtime implements them; \citet{anthropic-tooluse} for the exact request and result
representation a provider emits and expects; and \demoref{dispatch.py} and \demoref{workspace.py}
for the checks between a request and its result.

\paragraph{Looking ahead.} In the execution model described here, a routing function selects a
node and the scheduler runs it in the next superstep. The next section keeps the same machine
and changes this scheduling mechanism: a computation emits an event, which the runtime dispatches
to a registered handler. We then examine which component invokes the handler and when. The tool
contract and dispatcher remain the same. The new scheduling mechanism requires corresponding
evidence that each transition occurred once.
